\section{Finite Jacobi readouts and Gaussian quadrature}
\label{sec:readouts}

We now fix a time \(\lambda\), and suppress it in \(\mu,J,P_n\).
The principal application below is \(\lambda=0\). Write
\begin{equation}
 h_n^*=\int P_n^2\,\dd\mu=\frac{h_n}{Z_\lambda},\qquad
 q_n=\frac{P_n}{\sqrt{h_n^*}},\qquad
 h_0^*=1,\qquad h_n^*=\prod_{j=1}^n\beta_j.
 \label{eq:probabilitynorm}
\end{equation}
Thus \(q_0=1\). Let \(e_n\) denote the sequence with a one at site
\(n\) and zeros elsewhere. In this basis, multiplication by \(u\)
is the self-adjoint Jacobi operator of Section~\ref{sec:evolution}.
Its first finite block and the associated scalar observation are
\begin{equation}
 J_N=\begin{pmatrix}
 0&a_1&&&\\a_1&0&a_2&&\\&a_2&0&\ddots&\\
 &&\ddots&\ddots&a_{N-1}\\&&&a_{N-1}&0
 \end{pmatrix},\qquad
 G_N(z)=e_0^{\mathsf T}\e^{\ii zJ_N}e_0.
 \label{eq:finite-readout}
\end{equation}
Here the displayed pattern means the \(N\)-site matrix on
\(0,\ldots,N-1\), with \(J_1=(0)\). A matrix exponential is defined
by its convergent power series. The two appearances of \(e_0\) select
the entry describing departure from, and return to, the constant
polynomial. We call this entry a \emph{readout}.

\subsection{Nodes, weights and exact moments}

\begin{proposition}[Finite spectral measure]
\label{prop:gauss}
The eigenvalues \(x_{1,N},\ldots,x_{N,N}\) of \(J_N\) are real,
simple, and equal to the roots of \(P_N\). There are positive weights
\(\omega_{j,N}\), summing to one, such that
\begin{equation}
 Q_N f=\sum_{j=1}^N\omega_{j,N}f(x_{j,N}),\qquad
 G_N(z)=Q_N(\e^{\ii z\cdot}).
 \label{eq:gauss-readout}
\end{equation}
The rule \(Q_N\) integrates every polynomial of degree at most
\(2N-1\) exactly against \(\mu\). Moreover,
\begin{equation}
 \int u^{2N}\,\dd\mu-Q_N(u^{2N})=h_N^*,\qquad
 F^{(2N)}(0)-G_N^{(2N)}(0)=(-1)^Nh_N^*.
 \label{eq:first-walk}
\end{equation}
The node measure, and consequently \(G_N\), are even.
\end{proposition}

\begin{proof}
Expansion of \(\det(uI-J_N)\) at the last row gives the monic
recurrence \eqref{eq:recurrence}, with initial polynomials \(1,u\).
It therefore equals \(P_N(u)\). Real symmetry gives an orthonormal
eigenbasis. The row recurrence shows that an eigenvector whose first
entry is zero must be zero; an eigenspace cannot have dimension two,
since a linear combination would have first entry zero. Hence the
eigenvalues are simple and the squared first entries of the normalized
eigenvectors give strictly positive weights. Their sum is
\(\|e_0\|^2=1\). The finite spectral decomposition proves
\eqref{eq:gauss-readout}.

The moment \(\int u^k\,\dd\mu=\langle e_0,J^ke_0\rangle\)
is a sum of products of edges along length-\(k\) paths from site zero
back to itself. The finite matrix retains exactly the paths inside
the first \(N\) sites. A path reaching the first omitted site \(N\)
and returning needs at least \(2N\) steps. At that length there is
exactly one: \(0,1,\ldots,N,N-1,\ldots,0\). Its edge product is
\(a_1^2\cdots a_N^2=h_N^*\). This proves exactness and
\eqref{eq:first-walk}. Finally the diagonal matrix with entries
\((-1)^n\) conjugates \(J_N\) to \(-J_N\) and fixes \(e_0\),
proving symmetry of the weighted node measure.
\end{proof}

This is Gaussian quadrature: ``Gaussian'' refers to polynomial
exactness of the integration rule, not to a Gaussian source density.
The degree index \(N\) counts sites, not Taylor terms of \(G_N\);
the latter is generally an exponential sum, not a polynomial in \(z\).

\subsection{Complex convergence and the two zero problems}

\begin{proposition}[Convergence on compact complex sets]
\label{prop:compact}
For fixed \(\lambda\), the functions \(G_N\) converge to \(F\)
locally uniformly on \(\mathbb C\). More explicitly, with
\(m_{2k}=\int u^{2k}\,\dd\mu\),
\begin{equation}
 \sup_{|z|\le R}|F(z)-G_N(z)|
 \le\sum_{k=N}^{\infty}\frac{m_{2k}R^{2k}}{(2k)!}
 \longrightarrow0\qquad(R<\infty).
 \label{eq:compact-bound}
\end{equation}
\end{proposition}

\begin{proof}
Every return path has a nonnegative edge product, and the finite
matrix retains a subset of the full paths. Hence
\(0\le Q_N(u^{2k})\le m_{2k}\) for every \(k\). Lower moments agree and odd
moments vanish. The two exponential series therefore give
\eqref{eq:compact-bound}; its majorant sums to
\(\int\cosh(Ru)\,\dd\mu<\infty\), by \eqref{eq:tail}.
\end{proof}

There are two distinct zero questions:
\begin{equation}
 P_N(u)=\det(uI-J_N)=0
 \quad\hbox{and}\quad
 G_N(z)=\sum_j\omega_{j,N}\e^{\ii zx_{j,N}}=0.
 \label{eq:two-zero-problems}
\end{equation}
The first determines source-coordinate nodes. The second describes
cancellation in observation height. Classical Jacobi interlacing
concerns the first problem \cite[\S4.1, Theorem~4.5]{teschl}.
For the second, direct calculation at the first three degrees gives
\begin{equation}
 \begin{aligned}
 G_1(z)&=1,\qquad G_2(z)=\cos(\sqrt{\beta_1}\,z),\\
 G_3(z)&=\frac{\beta_2}{\beta_1+\beta_2}
       +\frac{\beta_1}{\beta_1+\beta_2}
          \cos\bigl(\sqrt{\beta_1+\beta_2}\,z\bigr).
 \end{aligned}
 \label{eq:small-heads}
\end{equation}
Indeed, \(J_3\) has nodes \(0,\pm\sqrt{\beta_1+\beta_2}\).
Symmetry, total mass one and second moment \(\beta_1\) determine
the weights in \eqref{eq:small-heads}.

For the unit-variance Gaussian source, \(\beta_1=1,\beta_2=2\).
Its exact transform is \(\e^{-z^2/2}\), but its three-site readout
has zeros
\begin{equation}
 G_3(z)=\tfrac23+\tfrac13\cos(\sqrt3z),\qquad
 z=\frac{(2k+1)\pi\pm\ii\arcosh 2}{\sqrt3},\quad k\in\mathbb Z.
 \label{eq:gaussian-counterexample}
\end{equation}
This elementary example excludes a generic inference from real Gauss
nodes to real Fourier-readout zeros. It does not test an additional
theta-specific hypothesis: Gaussian coefficients do not satisfy the
logarithmic envelope of Theorem~\ref{thm:theta}.

\subsection{Why polynomial phase is monotone}

The phase in this subsection counts zeros of the source polynomials,
not zeros of \(F\). It is background for a possible comparison method,
not a connecting theorem. In the Gaussian example above, the polynomials
have their usual real zeros although \(F(z)=\e^{-z^2/2}\) has none.
For a direct zero count of \(F\), the relevant phase is the change of
\(\arg F\) around a zero-free contour in the \(z\)-plane; the argument
principle counts the enclosed zeros. Section~\ref{sec:cluster-resolution}
uses precisely that contour information to follow unresolved clusters.

There is a concrete positive identity behind polynomial oscillation.
Subtract the recurrences at two real source points \(u,v\), multiply
by the opposite polynomial, and sum from degree zero to \(N-1\).
The interior terms cancel, leaving the Christoffel--Darboux identity
\begin{equation}
 (u-v)\sum_{j=0}^{N-1}q_j(u)q_j(v)
 =a_N\bigl(q_N(u)q_{N-1}(v)-q_{N-1}(u)q_N(v)\bigr).
 \label{eq:christoffel-darboux}
\end{equation}
Letting \(v\to u\) gives
\begin{equation}
 a_N\bigl(q_N'(u)q_{N-1}(u)-q_N(u)q_{N-1}'(u)\bigr)
       =\sum_{j=0}^{N-1}q_j(u)^2>0.
 \label{eq:polynomial-wronskian}
\end{equation}
The strict inequality uses \(q_0=1\). Consecutive polynomials
have no common zero, since the recurrence would propagate one
back to \(q_0\). Therefore their pair has a continuous unwrapped
phase on the real line. With the convention
\(\theta_N(u)=\arg(q_{N-1}(u)+\ii q_N(u))\), its derivative is
\begin{equation}
 \theta_N'(u)=
 \frac{\sum_{j=0}^{N-1}q_j(u)^2}
 {a_N\bigl(q_{N-1}(u)^2+q_N(u)^2\bigr)}>0.
 \label{eq:polynomial-phase}
\end{equation}
Here \emph{unwrapped} means choosing the angle continuously rather
than reducing it to an interval of length \(2\pi\).
Zeros of \(q_N\) occur at integer multiples of \(\pi\);
zeros of \(q_{N-1}\) occur at the intervening half-integer multiples.
This is an elementary phase explanation of their interlacing.
It concerns differentiation in \(u\), and it holds for every
positive Jacobi chain at nondegenerate degrees, including the
control sources. It supplies no corresponding positive identity
for differentiation of the Fourier readouts in \(z\).
